% Nouveau EGA V (cote 25): archive cote 25, site page 125 (PDF page 126).
% typescript. Draft compared: none.
% AI-transcribed by Claude Opus 5.5 (claude-opus-5-5, adaptive thinking, effort high), Message Batches API, batch ega25b (msgbatch_01A7n5YuXubUPNK1zsepi8RG), 2026-11-13; not yet checked by eye.
\scanpage{25/125}
\var{\inserted{compte tenu que l'hyp. sur $f$ implique $f_*(\underline{O}_X)\simeq\underline{O}_S$, d'où $f_*(\underline{O}_X^*)\simeq\underline{O}_S^*$}}{compte tenu que l'hyp. sur $f$ implique $f_*(\underline{O}_X)\simeq\underline{O}_S$, d'où $f_*(\underline{O}_X^*)\simeq\underline{O}_S^*$}{insertion}{ega25_25_125-v0001}
\mspage{\hand{28}}

par $\overset{\hand{\scriptscriptstyle\vee}}{P}$ , et associés à $\underline{L}$ (dans le sens de \hand{4}.1.) correspondent biunivo-%
quement aux\hand{*} homomorphismes \uline{surjectifs} $Q\hand{\longrightarrow}\overset{\hand{\scriptscriptstyle\vee}}{E}$\hand{,} modulo \var{\inserted{multiplication par des}}{multiplication par des}{insertion}{ega25_25_125-v0002} sections de $\underline{O}_S^{\hand{*}}$ .\hand{$\blacktriangledown$}
\var{\inserted{\var{\deleted{i.e. aux sections du faisceau $\mathcal{H}om_S(Q,\overset{\hand{\scriptscriptstyle\vee}}{E})^*/\underline{O}_S^*$, où le \unsure{premier} * \unsure{indique}}}{}{deletion}{ega25_25_125-v0004}}}{\var{\deleted{i.e. aux sections du faisceau $\mathcal{H}om_S(Q,\overset{\hand{\scriptscriptstyle\vee}}{E})^*/\underline{O}_S^*$, où le \unsure{premier} * \unsure{indique}}}{}{deletion}{ega25_25_125-v0004}}{insertion}{ega25_25_125-v0003}
\var{\margin{\unsure{Pour faisceaux d'ensembles} \illegible{2} \var{\deleted{\illegible{2}}}{}{deletion}{ega25_25_125-v0006} \unsure{préciser}}}{\placedmargin{\unsure{Pour faisceaux d'ensembles} \illegible{2} \var{\deleted{\illegible{2}}}{}{deletion}{ega25_25_125-v0006} \unsure{préciser}}}{marginal note placed}{ega25_25_125-v0005}
L'existence de $Q$ se ramène par le procédé brévété au cas S noethé-%
rien, et dans ce cas n'est autre que III 7.7.6. (l'hypothèse sur les
fibres de X étant d'ailleurs inutil\unsure{e|è}). Comme $\underline{E}$ est loc libre de type fini
la donnée d'un hom $\underline{E}\hand{\longrightarrow} f(\underline{L})$ équivaut à celle d'une section de
$\underline{L}\otimes_{\underline{O}_S}\overset{\hand{\scriptscriptstyle\vee}}{\underline{E}}$ , donc à celle d'un homomorphisme $Q\hand{\longrightarrow}\overset{\hand{\scriptscriptstyle\vee}}{\underline{E}}$ . Il reste à exprimer
que la condition envisagée dans \hand{4}.1. est bien vérifiée, ce qui (grâce à
l'hypothèse faite sur les fibres de X/S) se réduit à vérifier que
fibre par fibre l'homomorphisme correspondant $E(s)\hand{\longrightarrow}$ \var{\xout{Q(s)}}{}{typewriter strike-out}{ega25_25_125-v0007}\hand{$\vee$}
$H^0(X_s,L_s)\simeq\mathrm{Hom}_{k(s)}(Q(s),k(s))$ est injectif, ou encore $Q(s)\hand{\longrightarrow}\overset{\hand{\scriptscriptstyle\vee}}{E}(s)$ surjec%
tif, ce qui par Nakayama signffie aussi que $Q$ \quad $E$ est surjectif. Le ``mo
dulo sections de $\underline{O}_S$ '' provient du ``modulo isomorphisme'' dans \hand{4}.1.\hand{/}

On peut \var{\xout{xxxxxxxxx}}{}{typewriter strike-out}{ega25_25_125-v0008}\var{\inserted[typed]{interpréter}}{interpréter}{insertion}{ega25_25_125-v0009} \hand{4}.1. par une autre voie, en utilisant le fait
que $P(\underline{Q})$ représente le \var{\xout{xxxxxxxx}}{}{typewriter strike-out}{ega25_25_125-v0010} sous-foncteur de $\mathrm{Div}_{X/S}$ défini par
$\underline{L}$, en v ertu du n$^{\circ}$2. Par suite, un\var{\xout{x}}{}{typewriter strike-out}{ega25_25_125-v0011}système linéaire de diviseurs paramétr
par $\overset{\hand{\scriptscriptstyle\vee}}{P}$ \var{\inserted[typed]{et associé à $\underline{L}$}}{et associé à $\underline{L}$}{insertion}{ega25_25_125-v0012} s'interprète comme un morphisme $\overset{\hand{\scriptscriptstyle\vee}}{P} = P(\overset{\hand{\scriptscriptstyle\vee}}{\underline{E}})\hand{\longrightarrow} P(Q)$; le caractère
\uline{linéaire} de la famille de diviseurs définie par un tel morphisme
s'interprète alors par le fait que ce morphisme soit ``linéaire'', i.e.
précisément défini par un morphisme surjectif $Q\hand{\longrightarrow}\overset{\hand{\scriptscriptstyle\vee}}{\underline{E}}$ . \var{\margin{$\ulcorner$}}{\placedmargin{$\ulcorner$}}{marginal note placed}{ega25_25_125-v0013} On voit aussi
dans\var{\xout{xx}}{}{typewriter strike-out}{ega25_25_125-v0014} ce cas que le morphisme $\overset{\hand{\scriptscriptstyle\vee}}{P}\hand{\longrightarrow}\mathrm{Div}_{X/S}$ est un monomorphisme \var{\deleted{\illegible{2}}}{}{deletion}{ega25_25_125-v0015}
(puisque $\overset{\hand{\scriptscriptstyle\vee}}{P}\hand{\longrightarrow} P(Q)\hand{\hookrightarrow}\mathrm{Div}_{X/S}$ l'est), fait d'ailleurs \var{\inserted{plus}}{plus}{insertion}{ega25_25_125-v0016} général, cf corol-%
laire plus bas. Convenons \var{\inserted[typed]{alors}}{alors}{insertion}{ega25_25_125-v0017} de dire que d\unsure{e|a}ux familles linéaires de divi-%
seurs de X/S, paramétrées par les fibrés projectifs $\overset{\hand{\scriptscriptstyle\vee}}{P},\overset{\hand{\scriptscriptstyle\vee}}{P}'$ sont \uline{isomorph-}%
\uline{es} s'ils sont transformés l'un en l'autre par un S-isomorphisme $\overset{\hand{\scriptscriptstyle\vee}}{P}\hand{\longrightarrow}\overset{\hand{\scriptscriptstyle\vee}}{P}'$ ,
(lequel sera d'ailleurs unique grâce au fait qu'on a des monomorphismes
dans $\mathrm{Div}_{X/S}$ ). \var{\deleted{On appelera système linéaire de diviseurs dans X/S}}{}{deletion}{ega25_25_125-v0018}
(\var{\deleted{sans référence à un fibré projectif précisé qui le paramétriserait) une}}{}{deletion}{ega25_25_125-v0019}

\mspage{\hand{125}}
